
\documentclass{article}
\usepackage{colortbl}
\usepackage{makecell}
\usepackage{multirow}
\usepackage{supertabular}

\begin{document}

\newcounter{utterance}

\centering \large Interaction Transcript for game `eqbench', experiment `eq\_bench', episode 151 with reflexion{-}Qwen3.5{-}4B.
\vspace{24pt}

{ \footnotesize  \setcounter{utterance}{1}
\setlength{\tabcolsep}{0pt}
\begin{supertabular}{c@{$\;$}|p{.15\linewidth}@{}p{.15\linewidth}p{.15\linewidth}p{.15\linewidth}p{.15\linewidth}p{.15\linewidth}}
    \# & \multicolumn{2}{c}{Player} && \multicolumn{2}{c}{Game Master} \\
    \hline

    \theutterance \stepcounter{utterance}  
    & & & \multicolumn{4}{p{0.6\linewidth}}{
        \cellcolor[rgb]{0.9,0.9,0.9}{
            \makecell[{{p{\linewidth}}}]{
                \texttt{\tiny{[P1$\langle$GM]}}
                \\ 
\texttt{Your task is to predict the likely emotional responses of a character in this dialogue:} \\
\\ 
\texttt{Dr. Rigby: You remind me of someone I used to know.} \\
\texttt{Ellie: Oh, yeah? Some loser you used to pity?} \\
\texttt{Dr. Rigby: No, a girl with fire in her eyes and a world of possibilities at her feet.} \\
\texttt{Ellie: Sounds like a fairy tale. No one's got a world of possibilities.} \\
\texttt{Dr. Rigby: Not with that attitude, no. But imagine, just for a moment, a different reality.} \\
\texttt{Ellie: I don't have time for your riddles, lady.} \\
\texttt{Dr. Rigby: It's not a riddle, it's a wormhole.} \\
\texttt{Ellie: A what?} \\
\texttt{Dr. Rigby: A wormhole. A shortcut through space and time. What if you could reach out and change your destiny?} \\
\texttt{Ellie: You're nuts.} \\
\texttt{Dr. Rigby: Perhaps. But what if I told you that the future is not set in stone? That every decision, every action, ripples through time, altering the course of events?} \\
\texttt{Ellie: And what? I'm supposed to believe that my shitty decisions matter?} \\
\texttt{Dr. Rigby: More than you can imagine. You have the power to shape your own future.} \\
\texttt{Ellie: And how would you know?} \\
\texttt{{[}End dialogue{]}} \\
\\ 
\texttt{At the end of this dialogue, Dr. Rigby would feel...} \\
\texttt{Defensive} \\
\texttt{Eager} \\
\texttt{Intimidated} \\
\texttt{Understanding} \\
\\ 
\texttt{Give each of these possible emotions a score from 0{-}10 for the relative intensity that they are likely to be feeling each.} \\
\\ 
\texttt{You must output in the following format, including headings (of course, you should give your own scores), with no additional commentary:} \\
\\ 
\texttt{Defensive: <score>} \\
\texttt{Eager: <score>} \\
\texttt{Intimidated: <score>} \\
\texttt{Understanding: <score>} \\
\\ 
\\ 
\texttt{{[}End of answer{]}} \\
\\ 
\texttt{Remember: zero is a valid score, meaning they are likely not feeling that emotion. You must score at least one emotion > 0.} \\
\\ 
\texttt{Your answer:} \\
            }
        }
    }
     \\ \\

    \theutterance \stepcounter{utterance}  
    & \multicolumn{4}{p{0.6\linewidth}}{
        \cellcolor[rgb]{0.561,0.737,0.561}{
            \makecell[{{p{\linewidth}}}]{
                \texttt{\tiny{[P1$\rangle$GM]}}
                \texttt{Defensive: 0} \\
\texttt{Eager: 0} \\
\texttt{Intimidated: 0} \\
\texttt{Understanding: 0} \\
            }
        }
    }
    & & \\ \\

    \theutterance \stepcounter{utterance}  
    & & & \multicolumn{2}{p{0.3\linewidth}}{
        \cellcolor[rgb]{0.9,0.9,0.9}{
            \makecell[{{p{\linewidth}}}]{
                \texttt{\tiny{[GM$|$GM]}}
                \texttt{ParseError: At least one score in the response must be greater than 0,but the scores are: {[}'0', '0', '0', '0'{]}} \\
            }
        }
    }
    & & \\ \\

    \theutterance \stepcounter{utterance}  
    & & & \multicolumn{2}{p{0.3\linewidth}}{
        \cellcolor[rgb]{0.9,0.9,0.9}{
            \makecell[{{p{\linewidth}}}]{
                \texttt{\tiny{[GM$|$GM]}}
                \texttt{game\_result = ABORT} \\
            }
        }
    }
    & & \\ \\

    \theutterance \stepcounter{utterance}  
    & & \multicolumn{4}{p{0.6\linewidth}}{
        \cellcolor[rgb]{0.95,0.95,0.95}{
            \makecell[{{p{\linewidth}}}]{
                \texttt{\tiny{[P1$\leftrightarrow$P]}}
                \texttt{Post{-}episode Reflexion (not sent to GM):} \\
\texttt{New plan: I failed because I outputted all four scores as 0, which likely violated the rule "You must score at least one emotion > 0." I need to analyze the dialogue to determine Dr. Rigby's emotional state. Dr. Rigby is trying to inspire Ellie, but Ellie is hostile, dismissive, and calls him "nuts." Dr. Rigby's final line ("You have the power to shape your own future") is a direct response to Ellie's rejection, attempting to reframe her perspective. He is not intimidated (he's the one speaking), not eager (he's trying to convince, not just express excitement), and not defensive (he's not being attacked, he's the aggressor in the conversation). The most fitting emotion is "Understanding" in the sense of empathetic insight or perhaps a realization of the situation, but looking closer at the options: "Defensive" is unlikely. "Eager" is possible if he's desperate to change her mind. "Intimidated" is unlikely. "Understanding" fits the philosophical nature of his speech. However, the most logical interpretation of his reaction to being called "nuts" and having his "riddles" dismissed is that he feels a mix of frustration and a determined belief in his message. Let's re{-}evaluate "Defensive": He says "Not with that attitude, no" and "It's not a riddle," which are defensive justifications. But the prompt asks for his feeling *at the end*. At the end, he is asserting his point. "Understanding" seems the best fit for a character who sees the potential in the other despite their current state. I will assign a score of 8 to Understanding and 0 to the others, ensuring at least one is > 0.} \\
\\ 
\texttt{Defensive: 0} \\
\texttt{Eager: 0} \\
\texttt{Intimidated: 0} \\
\texttt{Understanding: 8} \\
            }
        }
    }
    & \\ \\

\end{supertabular}
}

\end{document}
